\documentclass[11pt]{article}
\usepackage[a4paper,margin=2.3cm]{geometry}
\usepackage{amsmath,amssymb,amsthm,booktabs,microtype}
\usepackage[hidelinks]{hyperref}
\newtheorem{proposition}{Proposition}
\newtheorem{definition}{Definition}
\theoremstyle{remark}\newtheorem{remark}{Remark}
\newcommand{\R}{\mathbb{R}}
\newcommand{\E}{\mathbb{E}}
\newcommand{\N}{\mathcal{N}}
\title{Why the mixture fit overfits, and the fix}
\author{Version-4 basket model, branch \texttt{v4-mixture}}
\date{28 October 2026}
\begin{document}
\maketitle

\noindent The mixture fit on the synthetic world \texttt{disc\_c48\_s3} gets worse on held-out trips after
3.5 epochs while it keeps getting better on training trips. This note (1) states what was measured,
(2) locates the parameters responsible, (3) explains why those parameters overfit and the
taste $z$ does not, and (4) derives the fix. Nothing in Sections~4--5 has been run yet.

%======================================================================
\section{Setting and notation}

The data are trips $t=1,\dots,T$. Trip $t$ belongs to household $h(t)$, and $S_t$ is its basket.
The training trips of household $h$ are $\mathcal T_h$, and $n_h=|\mathcal T_h|$ is their number.
In this world there are $T=7{,}695$ training trips and $H=928$ households, with $n_h$ of median 7 and mean 8.3.

The model (\texttt{docs/mixture\_theory.pdf}) has three kinds of unknowns.
\begin{enumerate}
\item \textbf{The taste of a trip}, $z_t\in\R^r$ ($r=16$). It has the prior $z_t\sim\N(0,I)$ and is
  \emph{integrated out}:
  \begin{equation}\label{eq:marg}
  p(S_t)=\int \N(z;0,I)\,p(S_t\mid z)\,dz .
  \end{equation}
\item \textbf{Household parameters}, one vector per household:
  \[\eta_h=(\theta_h,\gamma_h)\in\R^{d_h},\qquad d_h=32+8=40.\]
  Here $\theta_h\in\R^{32}$ is the household taste; its last coordinate is the household's size
  shift. $\gamma_h\in\R^8$ gives the household's price slopes. Both enter the utility $b_{tj}$ of
  every product $j$ on every trip of $h$.
\item \textbf{Product and shared parameters}, written together as $\omega$. They include the
  product context ($\alpha_j,\beta_j,\mu_j,\zeta_j$, display and mailer weights, recency loadings),
  about 57 numbers per product. A product is bought a median 50 times in training, and 10\% of
  products 6 times or fewer. They also include $\lambda_j$, $\phi_j$, $\psi_c$ and the size terms
  $\rho_0,\rho_c,f_c$.
\end{enumerate}
The fit maximises over $\eta=(\eta_h)_h$ and $\omega$:
\begin{equation}\label{eq:fit}
J(\eta,\omega)=\sum_{t}\log p_{\eta_{h(t)},\omega}(S_t)\;-\;\Pi(\omega)\;-\;\Pi_{\rm pool}(\eta,\omega).
\end{equation}
In \eqref{eq:fit}, $\Pi$ is the Gaussian prior on $\phi,\psi,\rho_c,f$, on the price terms (centred at the
base model) and on $u_c$. The term $\Pi_{\rm pool}$ is a pooling penalty on the products
$\alpha_j\cdot\theta_h$, with weight $0.5$ in total over all $T$ trips. In practice it does not
constrain $\theta_h$. So $\theta_h$ has, in effect, a flat prior.

%======================================================================
\section{What was measured}

All scores below are on the same 1{,}024 held-out trips and the same 1{,}024 training trips, with
$K=256$ importance draws per trip and the same random numbers for every checkpoint. The
differences are paired, trip by trip.

\begin{center}\small
\begin{tabular}{lrrrr}\toprule
checkpoint & held-out LL & training LL & training $-$ held-out & generated size (data 8.69)\\\midrule
step 242 (2 epochs) & $-27.29$ & $-26.76$ & 0.54 & 8.57\\
step 420 (best)     & $-26.84$ & $-25.84$ & 1.00 & 8.68\\
step 1{,}210 (final)& $-27.02$ & $-24.54$ & 2.48 & 7.92\\\bottomrule
\end{tabular}
\end{center}
From step 420 to 1{,}210, held-out trips lose $0.18\pm0.04$ nats and training trips gain
$1.31\pm0.04$.

\paragraph{Which parameters.} Take the final checkpoint and put one block back to its step-420
value. If a block's late changes only memorised the training trips, putting it back should lose
on training trips and gain on held-out trips.
\begin{center}\small
\begin{tabular}{lrrl}\toprule
block put back to step 420 & held-out & training & reading\\\midrule
household $\eta_h=(\theta_h,\gamma_h)$ & $+0.43\pm0.02$ & $-0.47\pm0.02$ & memorised\\
product context ($\alpha,\beta,\mu,\zeta,w,\dots$) & $+0.25\pm0.02$ & $-0.49\pm0.02$ & memorised\\
interactions ($\phi,\psi,\lambda$) & $-0.87\pm0.05$ & $-0.70\pm0.05$ & learned\\
size ($\rho_0,\rho_c,f$) & $-0.49\pm0.04$ & $-0.64\pm0.04$ & coupled to the rest\\\bottomrule
\end{tabular}
\end{center}
\begin{itemize}
\item The late changes of the household and product-context blocks only memorise the training trips.
\item The interactions keep learning something that also holds on held-out trips.
\item The size block has adjusted to the other blocks; put back on its own, it generates baskets of 11 products.
\item The final checkpoint with only the household block put back scores $-26.59$ on held-out trips,
  better than the best checkpoint ($-26.84$). So the best checkpoint is already partly overfitted.
\end{itemize}
Splitting the training gain by household gives the same picture. Households with 0--4 training trips
gain $+2.46$ per trip; those with 12 or more gain $+1.06$.

%======================================================================
\section{Why: the taste is integrated, the household parameters are not}

\subsection{\texorpdfstring{The taste $z_t$ cannot overfit}{The taste z cannot overfit}}
In \eqref{eq:marg}, $z_t$ is averaged over its prior; it is never fitted. A basket that would be
explained well by some extreme $z$ gets credit only in proportion to $\N(z;0,I)$ times the volume of
such $z$. There is no free per-trip number for the optimiser to tune. This is why the interaction
block keeps generalising.

\subsection{The household parameters are incidental parameters}
\begin{definition}
A parameter is \emph{incidental} when its number grows with the data, and each copy is informed by
a bounded number of observations. Here there are $H$ copies of $\eta_h$, and each is informed by $n_h$ trips.
\end{definition}
\noindent For a regular model fitted by maximum likelihood, the standard expansion of the fitted
log-likelihood (the expansion behind AIC and TIC) gives, for the $n_h$ trips of household $h$,
\begin{equation}\label{eq:optimism}
\E\Big[\sum_{t\in\mathcal T_h}\log p_{\hat\eta_h}(S_t)\Big]-\E\Big[\sum_{t\in\mathcal T_h}\log p_{\hat\eta_h}(S'_t)\Big]
\;\approx\;\operatorname{tr}\!\big(\mathcal I_h^{-1}\mathcal J_h\big)\;\approx\;d_h .
\end{equation}
Here $\hat\eta_h$ is the maximum-likelihood estimate from $\mathcal T_h$, and $S'_t$ are fresh trips
of the same household. $\mathcal I_h$ is the Fisher information of the $n_h$ trips and
$\mathcal J_h$ the variance of their score; the trace equals $d_h$ when the model is correct.

So, at convergence, the training-minus-held-out gap caused by the household block is about
$d_h/n_h$ nats per trip. With $d_h=40$ and $n_h=7$ that is about $5.7$ nats. It does \emph{not}
shrink as more households are added: $n_h$ stays at 7. The measured gap ($\approx0.9$ from this
block after 10 epochs) is on its way there. The bigger gain for households with few trips is the
$1/n_h$ in this bound. For products the same count gives about 57 numbers per product against a
median 50 purchases. That gap is smaller, as measured, but not negligible for rare products.

\subsection{Why the base model did not show it: early stopping was its prior}
The base model (\texttt{fit\_additive.py}) fits the same $\eta_h$ and product context with weight
decay $10^{-5}$, which is negligible. It stopped at its best validation score, iteration 800 of
30{,}000. Early stopping acts as a prior.
\begin{proposition}\label{prop:es}
Take a quadratic loss $\tfrac12(\eta-\hat\eta)^\top A(\eta-\hat\eta)$ with $A\succ0$. Gradient descent
with step $\epsilon$ from $\eta^{(0)}$, after $k$ steps, gives
$\eta^{(k)}-\hat\eta=(I-\epsilon A)^k(\eta^{(0)}-\hat\eta)$. Along each eigen-direction of $A$ with
eigenvalue $a$, this is the ridge (Gaussian-prior) estimate centred at $\eta^{(0)}$ with prior
precision $\tau$, to first order, when $(1-\epsilon a)^k=\tau/(a+\tau)$, that is, $\tau\approx 1/(\epsilon k)$ for $\epsilon a k\ll1$.
\end{proposition}
\begin{proof}
The ridge estimate is $\eta_\tau=\hat\eta+\tau(A+\tau I)^{-1}(\eta^{(0)}-\hat\eta)$. Along an
eigenvector with eigenvalue $a$ its shrinkage factor is $\tau/(a+\tau)$, and that of gradient descent
is $(1-\epsilon a)^k$. For $\epsilon a k\ll 1$, $(1-\epsilon a)^k\approx1-\epsilon ak$ and
$\tau/(a+\tau)\approx 1-a/\tau$; these match when $\tau=1/(\epsilon k)$.
\end{proof}
\noindent The fits use Adam, not plain gradient descent, so for them the proposition holds only
qualitatively: fewer steps act as a stronger prior centred at the start.
So the base model's household and product-context terms were held near their starting
values by a prior of precision about $1/(\epsilon k)$. Its strength was set by where training stopped.
The mixture fit trains the same terms for 1{,}210 more steps. That weakens this implicit prior
until nothing holds $\eta_h$ back. Choosing the best checkpoint on validation just re-introduces
early stopping, one stopping time for all blocks. The blocks need different strengths: the
interactions want to keep training, and the household terms want to stop.

%======================================================================
\section{The fix: a hierarchical prior with an estimated scale}

Make the household and product-context parameters random effects with an explicit prior, in the
same way as $z_t$.
\begin{equation}\label{eq:prior}
\theta_h\sim\N(0,\sigma_\theta^2 I),\qquad \gamma_h\sim\N(\bar\gamma,\sigma_\gamma^2 I),\qquad
\xi_j\sim\N(0,\sigma_\xi^2 I),
\end{equation}
for each product-context block $\xi\in\{\alpha,\beta,\mu,\zeta,w,\dots\}$.
The scales $\sigma^2$ are \emph{estimated} rather than set by hand, and the stopping time no longer
matters.

\paragraph{Step 1: the penalised fit (MAP).} For fixed scales, maximise
\begin{equation}\label{eq:map}
J_\sigma(\eta,\omega)=\sum_t\log p_{\eta_{h(t)},\omega}(S_t)-\sum_h\frac{\lVert\theta_h\rVert^2}{2\sigma_\theta^2}
-\sum_h\frac{\lVert\gamma_h-\bar\gamma\rVert^2}{2\sigma_\gamma^2}-\sum_{\xi}\sum_j\frac{\lVert\xi_j\rVert^2}{2\sigma_\xi^2}-\Pi(\omega).
\end{equation}
This replaces the flat prior on $\theta_h$, and the price prior's fixed scale on $\gamma_h$, with
learned ones. The gradient of the likelihood term is unchanged; each block only gets its own
Gaussian term.

\paragraph{Step 2: the scales (empirical Bayes).} The scale that maximises the marginal likelihood
of the data, with $\eta_h$ integrated under \eqref{eq:prior}, satisfies a fixed point. Under the
Laplace approximation to each household's posterior, the EM update is
\begin{equation}\label{eq:em}
\sigma_\theta^2\leftarrow\frac{1}{H\,d_\theta}\sum_h\Big(\lVert\hat\theta_h\rVert^2+\operatorname{tr}\big(A_h+\sigma_\theta^{-2}I\big)^{-1}\Big),
\end{equation}
where $A_h=-\nabla^2_{\theta_h}\sum_{t\in\mathcal T_h}\log p(S_t)$ is the curvature from the household's own trips.
\begin{itemize}
\item The trace term is the posterior variance that MAP leaves out. Without it, the update
  shrinks $\sigma^2$ too far toward zero when $n_h$ is small.
\item The same update holds for $\sigma_\gamma^2$ (with $\hat\gamma_h-\bar\gamma$, and $\bar\gamma$
  set to the mean of $\hat\gamma_h$) and for each $\sigma_\xi^2$ (over products).
\item $A_h$ can be replaced by its expected form, the empirical Fisher matrix
  $\sum_{t\in\mathcal T_h}g_tg_t^\top$, with $g_t=\nabla_{\theta_h}\log p(S_t)$. Each $g_t$ is a
  by-product of the training step.
\end{itemize}
The update is run every few hundred steps, like the whitening step.

\begin{proposition}\label{prop:shrink}
At a fixed point of \eqref{eq:map}--\eqref{eq:em}, a household with few trips is pulled toward the
population: along an eigen-direction of $A_h$ with eigenvalue $a$, its estimate is the
maximum-likelihood one shrunk by $a/(a+\sigma_\theta^{-2})$. Since $a$ grows with $n_h$, the shrinkage
fades as the household gets more trips. The optimism \eqref{eq:optimism} becomes, approximately, the effective
number of parameters, $\sum_i a_i/(a_i+\sigma_\theta^{-2})<d_h$, instead of $d_h$.
\end{proposition}
\begin{proof}
This is the quadratic case of \eqref{eq:map}: the penalised optimum is
$(A_h+\sigma^{-2}I)^{-1}A_h\hat\eta_h^{\rm ML}$. The effective number of parameters of a
ridge-penalised fit is $\operatorname{tr}\big((A_h+\sigma^{-2}I)^{-1}A_h\big)$; in \eqref{eq:optimism}
the inverse information is replaced by the inverse penalised curvature, which gives this trace
when the model is correct.
\end{proof}

\paragraph{What does not change.} The taste $z_t$, its prior $\N(0,I)$, the likelihood estimate,
generation, whitening and moment matching are all unchanged. The interactions keep training freely,
under their existing prior.

%======================================================================
\section{The complete version: integrate the household parameters too}

Proposition~\ref{prop:shrink} is still a point estimate. The consistent treatment does for $\theta_h$
what \eqref{eq:marg} does for $z_t$: the probability of all of a household's baskets is
\begin{equation}\label{eq:hh}
p\big((S_t)_{t\in\mathcal T_h}\big)=\int \N(\theta;0,\sigma_\theta^2 I)\prod_{t\in\mathcal T_h}\Big[\int\N(z_t;0,I)\,p(S_t\mid z_t,\theta)\,dz_t\Big]d\theta .
\end{equation}
\begin{proposition}\label{prop:joint}
Hold $\gamma_h$ and $\omega$ fixed, and treat the within-category competition shift and the
history centring as fixed offsets at their current values. Then the joint posterior of $(\theta_h,z_{t})_{t\in\mathcal T_h}$
given the household's baskets is strongly log-concave, so it has a single peak.
\end{proposition}
\begin{proof}
With the offsets fixed, the utility of product $j$ on trip $t$ is affine in $(\theta,z_t)$:
$u_{tj}=c_{tj}+\theta\cdot a_j+z_t\cdot\phi_j$. Here $a_j$ collects the size shift, $\alpha_j$ and
$u_{c(j)}$. The presence term is $-z_t\cdot\psi_c$ per extra product of category $c$, also affine
in $z_t$. Hence the energy $E_{z_t}(S,\theta)$ of every basket $S$ is affine in $(\theta,z_t)$.
So $\log p(S_t\mid z_t,\theta)=E_{z_t}(S_t,\theta)-\log Z_{z_t}(\theta)$, and
$\log Z_{z_t}(\theta)=\log\sum_S e^{E_{z_t}(S,\theta)}$ is the log of a sum of exponentials of affine
functions, hence convex in $(\theta,z_t)$ jointly. The sum over trips of concave terms, plus the strictly concave log-prior, is strongly concave.
\end{proof}
\noindent The price slopes $\gamma_h$ enter through a softplus, which breaks affinity, so they stay as in
Section~4. This version removes the household optimism entirely, but it costs much more: every step
needs the joint peak over a whole household's trips, and batches must be formed by household.
Section~4 is the first step; this version is the one to move to if Section~4's held-out gain stays
well short of what Section~2 shows is available.

%======================================================================
\section{What the fix must show (stated before any run)}
\begin{enumerate}
\item Held-out log-likelihood no longer falls as training continues: the last checkpoint is within
  noise of the best.
\item The training-minus-held-out gap stays near the value at step 242 (about 0.5), instead of reaching 2.5.
\item The block swap of Section~2 for the household and product-context blocks changes the held-out
  score by no more than its standard error.
\item Held-out log-likelihood at least $-26.59$, the score of the final checkpoint with only the
  household block put back.
\item Generated basket size within 0.2 of the data (8.69), with the 99\% quantile near 22.8.
\item The competition estimate stays near the truth, 0.05.
\end{enumerate}

\section{Open points}
\begin{itemize}
\item The real-data mixture fit also dropped after step 240 (by 0.66), yet its household and context
  terms were frozen there. This note does not explain that drop. The same block swap on the real
  checkpoints (interactions against size terms) is the next measurement there.
\item Update \eqref{eq:em} uses the Laplace approximation. With $n_h=1$ or 2 a household's posterior
  is far from Gaussian, and the trace term is then an approximation.
\end{itemize}

\end{document}
